\begin{pastexam}[label={px:nccu-107-dsa-06}]{政大資科所 107 學年度 第 6 題}
（問答）

(10\%) Design a polynomial-time algorithm to compute the maximum spanning tree, i.e., the spanning
tree that has the largest total edge weight.
\end{pastexam}

\begin{pastsoln}
\textbf{解讀}。 圖是連通的無向加權圖，$n$ 個頂點、$m$ 條邊（不連通就沒有生成樹，下面的做法會得到
每個元件的最大生成樹，也就是最大生成森林）。 權重可以是任何實數。

\textbf{想法：權重取負號}。 令 $w'(e) = -w(e)$。 對任何生成樹 $T$，$w'(T) = -w(T)$，所以 $T$ 讓 $w$
的總和最大，若且唯若它讓 $w'$ 的總和最小。 最大生成樹就是 $w'$ 下的最小生成樹。 Kruskal 演算法
（資結書第~24.2~節）的正確性只用到切割性質，不要求權重非負，所以可以直接對 $w'$ 做。 對 $w'$ 由小到大
排序，就是對 $w$ 由大到小排序，於是演算法就是「Kruskal，但邊由重到輕看」：
\begin{pseudo}[unbreakable]{最大生成樹}
\Function{Max-Spanning-Tree}{$G, w$}
    \State $T \gets \emptyset$
    \For{每個頂點 $v$}
        \State \Call{Make-Set}{$v$}
    \EndFor
    \State 把所有邊依 $w$ \textbf{由大到小}排序
    \For{依序每條邊 $\{u, v\}$}
        \If{\Call{Find-Set}{$u$} $\ne$ \Call{Find-Set}{$v$}} \Comment{不會成環}
            \State $T \gets T \cup \{\{u, v\}\}$
            \State \Call{Union}{$u, v$}
        \EndIf
    \EndFor
    \State \Return $T$
\EndFunction
\end{pseudo}

\textbf{正確性}。 上面已經說明：這就是在 $w' = -w$ 上跑 Kruskal（資結書定理~24.6），得到 $w'$ 的最小
生成樹，也就是 $w$ 的最大生成樹。 也可以直接用「反過來的切割性質」說明：對任何切割，跨過切割最\textbf{重}
的邊一定在某一棵最大生成樹裡（把資結書定理~24.1 套到 $w'$ 上）。 Kruskal 加入邊 $\{u, v\}$ 時，取 $u$ 所在
的元件當切割的一邊，所有比它重的邊都已經看過，跨過這個切割的邊裡它最重，所以加入它是安全的。

\textbf{例子}。 三角形 A--B、B--C、C--A 權重 $1, 2, 3$：先拿 C--A（$3$），再拿 B--C（$2$），A--B 會
成環，跳過。 最大生成樹的權重 $5$；另外兩棵生成樹是 $1 + 3 = 4$、$1 + 2 = 3$。

\textbf{時間}。 排序 $O(m \lg m)$；$m$ 次 \proc{Find-Set}、$n - 1$ 次 \proc{Union} 用互斥集合森林
（按秩合併加路徑壓縮，資結書第~13~章）幾乎是線性。 總共 $O(m \lg m) = O(m \lg n)$，是多項式時間。

另一種寫法：Prim 演算法每次改成挑「最重」的跨越邊，也是對的，理由一樣。
\end{pastsoln}
